\section{Essential incompleteness of arithmetic}\label{sec:arithmetic}

\subsection{Comparison in the fixed syntactic category}
\label{subsec:arithmetic-internal-comparison}

Fix the arithmetic signature and its realization in
$\mathcal B_{\PA}=\Tot(\Split(S_{\PA}))$ from
Theorem~\ref{thm:syntactic-generated-theory}.
Let $S=J_{\PA}(\operatorname{Sent})\in\mathcal B_{\PA}$ be
the sentence object, and let
$\nu:S\to S$ be syntactic negation.
We regard an object $X$ of $S_{\PA}$ as its identity splitting
$(X,1_X)$ when it occurs in $\mathcal B_{\PA}$.

Let $\Gamma$ be a computably enumerable set of arithmetic
sentences. Choose a $\Sigma_1$ formula
$\operatorname{Ax}_\Gamma(z)$ whose solutions in $\N$ are their codes.
The predicate
$\operatorname{Sent}(z)\land\operatorname{Ax}_\Gamma(z)$
defines an axiom subobject $A_\Gamma\hookrightarrow S$ in
$\mathcal B_{\PA}$. Write
\[
 P_\Gamma=P_{A_\Gamma,S},\qquad
 c_\Gamma=c_{A_\Gamma,S}:P_\Gamma\to S,\qquad
 t_\Gamma=\ran(c_\Gamma):S\rightharpoonup S.
\]
Here $c_\Gamma$ is the pullback of the conclusion morphism for all
formulas along $S\hookrightarrow\Fm$, and $t_\Gamma$ is its range in
$\Split(S_{\PA})$.
Choose an index $e_{\Gamma,S}$ representing $c_\Gamma$
as in Theorem~\ref{thm:syntactic-generated-theory}, and put
$\operatorname{Prov}_\Gamma(z)\equiv\exists p\,G_{e_{\Gamma,S}}(p,z)$.
This predicate is $\Sigma_1$ and implies
$\operatorname{Sent}(z)$ in PA. The morphism on $\omega$
representing $t_\Gamma$ is $[\operatorname{Prov}_\Gamma]_{\PA}$.
The identity of $S$ is $[\operatorname{Sent}]_{\PA}$.

Under the standard interpretation, the set
$G_{\N}P_\Gamma$ consists of the codes of finite derivations from $\Gamma$
whose conclusions are sentences. Consequently
\[
 \mathsf{Th}_\Gamma
   =\{s\in G_{\N}S:\N\models\operatorname{Prov}_\Gamma(s)\}
\]
is the set of sentence codes derivable from $\Gamma$.
Its refutation set is
$\mathsf{Ref}_\Gamma=(G_{\N}\nu)^{-1}(\mathsf{Th}_\Gamma)$.
The partial identities on these sets are
$\widehat F_{\N}(t_\Gamma)$ and
$(\widehat F_{\N}\nu)^*\widehat F_{\N}(t_\Gamma)$,
respectively.

Use the Turing morphism
$\tau_{\PA}=[u]_{\PA}:\omega\times\omega\rightharpoonup\omega$
and the coproduct isomorphism
$b:\omega\to\omega+\omega$ given by the even and odd codings.
In numerical coordinates the graph of $b$ is $y=x$; its standard
interpretation sends $2m$ to the first copy of $m$ and $2m+1$
to the second. Put
\[
 h=b\tau_{\PA}\Delta_\omega:\omega\rightharpoonup\omega+\omega,
 \qquad
 a_0=h^*(1_\omega+0_\omega),\qquad
 a_1=h^*(0_\omega+1_\omega).
\]
Thus $a_0,a_1:\omega\rightharpoonup\omega$ are restriction
idempotents in $S_{\PA}$.
Choose an index $e_d$ representing
$\tau_{\PA}\Delta_\omega$ and define
\[
 \alpha(x)\equiv\exists m\,G_{e_d}(x,2m),\qquad
 \beta(x)\equiv\exists m\,G_{e_d}(x,2m+1).
\]
By the PA-provable coproduct coding equations, the inverse images
of the two summands under $h$ have predicates $\alpha$ and $\beta$.
Thus $a_0=[\alpha]_{\PA}$ and $a_1=[\beta]_{\PA}$ in $S_{\PA}$.
The partial identities $F_{\N}(a_0),F_{\N}(a_1)$ have domains
$D_0=\{n:\N\models\alpha(n)\}$ and
$D_1=\{n:\N\models\beta(n)\}$.
Let $\Theta:\N\to\N$ be the primitive-recursive function
$\Theta(n)=\ulcorner\alpha(\bar n)\urcorner$, and choose a
primitive-recursive extension $\operatorname{Neg}:\N\to\N$
of negation on formula codes. On a formula code $s$,
$\operatorname{Neg}(s)$ adjoins the fixed unary negation constructor
at the root. The values of $\Theta$ belong to the sentence-code
subset $G_{\N}S$.
In arithmetic formulas,
these functions abbreviate their PA-provably total graph definitions.

\Needspace{11\baselineskip}
\begin{lemma}[Representation of the diagonal restriction idempotents]
\label{lem:arithmetic-diagonal-representation}
There is a total morphism
$\theta:\omega\to S$ in $\mathcal B_{\PA}$ such that
$(G_{\N}\theta)(n)=\ulcorner\alpha(\bar n)\urcorner$.
Suppose that every axiom of PA is derivable from the c.e.\
sentence set $\Gamma$. For every $M\equiv\N$, the following
inequalities hold in $\Pfn$:
\begin{align*}
 F_M(a_0)&\le
   (\widehat F_M\theta)^*\widehat F_M(t_\Gamma),\\
 F_M(a_1)&\le
   (\widehat F_M\theta)^*(\widehat F_M\nu)^*
                   \widehat F_M(t_\Gamma).
\end{align*}
\end{lemma}
\begin{proof}
Numeral construction and substitution in the fixed formula
$\alpha(x)$ compute $\Theta$ primitive-recursively.
PA proves that the numeral is a closed term, so it is admissible
for substitution, and that the resulting formula is a sentence.
The substitution realization of
Proposition~\ref{prop:st-substitution-realization} therefore gives
the stated total morphism $\theta$.

PA-provable functionality of $G$ and disjointness of even and odd
numbers give $\PA\vdash\forall x\,\neg(\alpha(x)\land\beta(x))$.
If $n\in D_0$, a finite computation trace ending in an even
output can be verified by a finite PA-proof. Introducing its
existential witnesses gives $\PA\vdash\alpha(\bar n)$.
If $n\in D_1$, the same argument gives a PA-proof of
$\beta(\bar n)$; the disjointness theorem then gives
a PA-proof of $\neg\alpha(\bar n)$.
Replacing each PA-axiom in either proof by its derivation from
$\Gamma$ gives a proof from $\Gamma$.
By the description of $\operatorname{Prov}_\Gamma$, the
universal closures of
\[
 \alpha(x)\longrightarrow\operatorname{Prov}_\Gamma(\Theta(x)),
 \qquad
 \beta(x)\longrightarrow
       \operatorname{Prov}_\Gamma(\operatorname{Neg}(\Theta(x)))
\]
are true in $\N$.
These two sentences hold in every $M\equiv\N$.
Interpreting their predicates gives the two inequalities, because $\widehat F_M$ preserves the range of
$c_\Gamma$ and inverse images.
\end{proof}

\begin{lemma}[A lift constructed by proof search]
\label{lem:arithmetic-proof-search}
For every c.e.\ sentence set $\Gamma$, there is a morphism
$\ell_\Gamma:\omega\rightharpoonup\omega+\omega$ in
$S_{\PA}$ with the following property. In any $M\models\PA$
for which
\[
 C(x)=\operatorname{Prov}_\Gamma(\Theta(x)),\qquad
 C'(x)=\operatorname{Prov}_\Gamma(\operatorname{Neg}(\Theta(x)))
\]
have disjoint truth sets, the partial function
\[
 s_\Gamma=\chi_{\omega,\omega}^{-1}F_M(\ell_\Gamma):
       F_M\omega\rightharpoonup F_M\omega+F_M\omega
\]
in $\Pfn$ satisfies
\[
 s_\Gamma^*(1_{F_M\omega}+0_{F_M\omega})=[C'(M)],\qquad
 s_\Gamma^*(0_{F_M\omega}+1_{F_M\omega})=[C(M)].
\]
Here $[C(M)]$ denotes the partial identity on
$\{x\in M:M\models C(x)\}$, and $[C'(M)]$ is defined in the
same way.
\end{lemma}
\begin{proof}
The formula
\[
 R(x,z)\equiv
   (C(x)\land z=2x+1)\lor(C'(x)\land z=2x)
\]
is $\Sigma_1$. Lemma~\ref{lem:syntactic-functional-search}
gives an index $a_R$ for which PA proves
\[
 G_{a_R}(x,z)\longrightarrow R(x,z),\qquad
 \exists z\,G_{a_R}(x,z)\longleftrightarrow C(x)\lor C'(x).
\]
Set $\ell_\Gamma=[a_R]_{\PA}:\omega\rightharpoonup\omega+\omega$.

In a model where $C,C'$ are disjoint, $R$ is functional. If
$M\models R(x,z)$, the second equivalence supplies
an output $z'$ of the interpreted program. The first implication
gives $M\models R(x,z')$, so $z'=z$. Thus the interpreted graph
is $R$. Decoding the even and odd outputs gives the
two inverse-image equations.
\end{proof}

\begin{theorem}[Incompleteness of an internal arithmetic theory]
\label{thm:arithmetic-internal-comparison}
Let $\Gamma$ be a consistent c.e.\ set of arithmetic sentences
from which every PA-axiom is derivable. Its sentence theory
$t_\Gamma$ in $\Split(S_{\PA})$ is incomplete.
Its interpretations under $\widehat F_{\N}$ and under every
$\widehat F_M$, where $M=\N^I/\mathcal U$ for a proper ultrafilter
$\mathcal U$ on a nonempty set $I$, are incomplete.
For the standard interpretation, there is a total morphism
$g:\omega\to\omega$ in $S_{\PA}$ such that, writing
$\gamma=F_{\N}(g):\N\to\N$, every sentence
$\alpha(\overline{\gamma(n)})$ is undecidable from $\Gamma$.
\end{theorem}
\begin{proof}
By first-order completeness there is a model $K\models\Gamma$.
The sentences of $\Gamma$ are interpreted in $K$. Their codes and
the codes of their derivations are interpreted in $\N$ or in an ultrapower
of $\N$.
Proposition~\ref{prop:standard-code-soundness} supplies the
soundness factorization for $c_\Gamma$ under $G_{\N}$.
Its validity subobject is disjoint from its negated inverse
image, so Proposition~\ref{prop:interpreted-theory} gives
\[
 \widehat F_{\N}(t_\Gamma)
    (\widehat F_{\N}\nu)^*\widehat F_{\N}(t_\Gamma)=0.
\]
For an ultrapower $M=\N^I/\mathcal U$, the same conclusion
follows from Theorem~\ref{thm:ultrapower-soundness}.

Apply Theorem~\ref{thm:logical-comparison} with
$\mathscr R=S_{\PA}$, $\mathcal A=\Pfn$ and $F=F_M$.
The target is split, its inverse images along total morphisms preserve binary
joins, and $F_M(1_\omega)\ne0$ because $M$ is nonempty.
Lemma~\ref{lem:arithmetic-diagonal-representation} gives the
two reduction inequalities. Lemma~\ref{lem:arithmetic-proof-search}
gives the partial conditional and its source lift.
Theorem~\ref{thm:logical-comparison} shows that the interpreted theorem domain
and its inverse image under interpreted negation have join strictly below the
identity. In the source it gives $t_\Gamma\vee\nu^*t_\Gamma\ne1_S$.
For $M=\N$, it supplies a total morphism $g:\omega\to\omega$
in $S_{\PA}$. With $\gamma=F_{\N}(g)$, the two inverse-image
equalities in that theorem become
\[
 \bigl((G_{\N}\theta)\gamma\bigr)^{-1}(\mathsf{Th}_\Gamma)
 =\varnothing,
 \qquad
 \bigl((G_{\N}\theta)\gamma\bigr)^{-1}(\mathsf{Ref}_\Gamma)
 =\varnothing.
\]
Thus neither $\alpha(\overline{\gamma(n)})$ nor its negation
is derivable from $\Gamma$, for any $n\in\N$.
\end{proof}

Di Paola and Heller list recursive inseparability of the theorem and
refutation sets as a consequence of definability of every recursive set
in a consistent arithmetical theory \cite[p.~594, (3)]{DPH}.
For a c.e.\ extension of PA, the function $\Theta$ sends $D_0$ into
its theorem set and $D_1$ into its refutation set by
Lemma~\ref{lem:arithmetic-diagonal-representation}.

\begin{corollary}[Essential incompleteness of arithmetic]
\label{cor:arithmetic}
\label{thm:arithmetic-comparison-inseparability}
For every consistent c.e.\ extension of PA, its theorem and
refutation sets of sentence codes are recursively inseparable.
Consequently every consistent c.e.\ extension of it is incomplete.
\end{corollary}
\begin{proof}
Choose a c.e.\ sentence presentation $\Gamma$ of the extension.
Suppose a recursive set $R\subseteq\N$ contains
$\mathsf{Th}_\Gamma$ and is disjoint from
$\mathsf{Ref}_\Gamma$. The recursive set
$C=\{n\in\N:\Theta(n)\in R\}$
contains $D_0$ and is disjoint from $D_1$, by
Lemma~\ref{lem:arithmetic-diagonal-representation}.
Let $P,N$ be the complementary partial identities on $C$ and
$\N\setminus C$. The total recursive function that sends
$n\in C$ to the second copy of $n$ and $n\notin C$ to the
first is their conditional in $\Pfn$.
Choose an index for this function, using the fixed numerical
coproduct coding. Its class is a morphism
$\ell:\omega\rightharpoonup\omega+\omega$ of $S_{\PA}$,
and its standard interpretation is this conditional.
Corollary~\ref{cor:separator}, applied to $F_{\N}$ with the
partial lift $\ell$, gives a contradiction.

Let $\Delta$ present a consistent c.e.\ extension of the theory of
$\Gamma$. Then $\mathsf{Th}_\Gamma\subseteq\mathsf{Th}_\Delta$ and
$\mathsf{Ref}_\Gamma\subseteq\mathsf{Ref}_\Delta$. If $\Delta$ were
complete, its two disjoint c.e.\ sets would cover the recursive set of
sentence codes. Parallel enumeration would decide which set contains
each sentence code. Assigning false to nonsentence codes would make
$\mathsf{Th}_\Delta$ a recursive separator of
$\mathsf{Th}_\Gamma$ and $\mathsf{Ref}_\Gamma$.
\end{proof}

\subsection{An indexed family of arithmetic theories}
\label{subsec:indexed-arithmetic-theories}

We construct one derivation object for all computably enumerable
axiom sets containing the chosen sentence axioms of PA. Put
$\mathcal B=\Tot(\Split(S_{\PA}))$, and let
$I:S_{\PA}\to\Split(S_{\PA})$ be the inclusion
$I(X)=(X,1_X)$ on objects and $I(f)=f$ for
$f:X\rightharpoonup Y$ in $S_{\PA}$.
Use the syntax realized in $\mathcal B$ by
Theorem~\ref{thm:syntactic-generated-theory}. In particular,
$S,\Fm\in\mathcal B$, and $j:S\hookrightarrow\Fm$ is the
sentence inclusion. All products, coproducts, and pullbacks in
the family construction below are taken in $\mathcal B$.
The parameter object is $E=I(\omega)=(\omega,1_\omega)$.
The interpretation $G_{\N}:\mathcal B\to\mathbf{Set}$ of
Section~\ref{sec:models} satisfies $G_{\N}E=\N$.

Choose a $\Sigma_1$ formula $\operatorname{Ax}_{\PA}(z)$
for the codes of the chosen sentence axioms of PA. Define
\[
 W(e,z)\equiv\exists y\,G(e,z,y),\qquad
 \operatorname{Ax}(e,z)\equiv
 \operatorname{Sent}(z)\land
       \bigl(\operatorname{Ax}_{\PA}(z)\lor W(e,z)\bigr).
\]
As $e$ varies over $\N$, the sets
$\{z\in\N:\N\models W(e,z)\}$ are precisely the c.e.\ subsets of
$\N$. Let $\Gamma_e$ be the set of sentences whose codes
satisfy $\operatorname{Ax}(e,z)$ in $\N$.
For each $e\in\N$, use the formula $\operatorname{Ax}(\bar e,z)$,
with $\bar e$ the numeral for $e$, to present the axiom subobject
$A_{\Gamma_e}\hookrightarrow S$. The notation
$\operatorname{Prov}_{\Gamma_e}$ below refers to this presentation.
Two enumerations with the same standard axiom codes need not define
provably equivalent predicates in $\PA$.

Let $J=J_{\PA}:\mathcal P_{\PA}\to\mathcal B$ be the
equivalence of Lemma~\ref{lem:syntactic-total-category}.
Put $\operatorname{Top}(x)\equiv(x=x)$, so that
$J(\operatorname{Top})=E$, and use the object
$S=J(\operatorname{Sent})$.
Let $\pi_0^{\mathrm{num}},\pi_1^{\mathrm{num}}:\N\to\N$
be the projections of the numerical pairing, and put
\[
 D(w)\equiv\operatorname{Sent}(\pi_1^{\mathrm{num}}w),\qquad
 A_0(w)\equiv\operatorname{Ax}
            (\pi_0^{\mathrm{num}}w,\pi_1^{\mathrm{num}}w).
\]
Let $r_0:D\to\operatorname{Top}$ and
$r_1:D\to\operatorname{Sent}$ in $\mathcal P_{\PA}$ be
represented, respectively, by
$D(w)\land e=\pi_0^{\mathrm{num}}w$ and
$D(w)\land z=\pi_1^{\mathrm{num}}w$.
The pairing identities give the product isomorphism
$v=\langle J(r_0),J(r_1)\rangle:J(D)\xrightarrow{\cong}E\times S$.
The predicate inclusion $a:A_0\hookrightarrow D$ in
$\mathcal P_{\PA}$ has graph $A_0(w)\land w'=w$.
Define the axiom subobject for the family in $\mathcal B$ by
$A=J(A_0)$ and $i=vJ(a):A\hookrightarrow E\times S$.

Let $L,R_1,R_2\in\mathcal B$ and
\[
 \lambda:L\to\Fm,\qquad
 p_n:R_n\to\Fm^n,\qquad q_n:R_n\to\Fm\quad(n=1,2)
\]
be the realized logical-axiom and rule morphisms of
Appendix~\ref{app:first-order-derivations}.
Put $F_E=E\times\Fm$ and define
\[
 R^E_0=(E\times L)+A,\qquad
 R^E_n=E\times R_n\quad(n=1,2).
\]
The nullary premise morphism is $p^E_0:R^E_0\to1$, the
unique morphism to the terminal object. Its conclusion morphism is
\[
 q^E_0=[1_E\times\lambda,(1_E\times j)i]:R^E_0\to F_E.
\]
For $n=1,2$, write
$u_n:E\times\Fm^n\to E$ and
$v_n:E\times\Fm^n\to\Fm^n$ for the product projections,
and $\pi_{n,k}:\Fm^n\to\Fm$ for the $k$th projection,
where $1\le k\le n$.
Define $d_n:E\times\Fm^n\to F_E^n$ by its components
$\langle u_n,\pi_{n,k}v_n\rangle:E\times\Fm^n\to F_E$,
and put
\[
 p^E_n=d_n(1_E\times p_n):R_n^E\to F_E^n,\qquad
 q^E_n=1_E\times q_n:R_n^E\to F_E.
\]
Every premise and the conclusion of a rule have the same
$E$-coordinate. These morphisms define the inference endofunctor
$H_E:\mathcal B/F_E\to\mathcal B/F_E$
of Appendix~\ref{app:closed-subobjects}.
Lemma~\ref{lem:syntactic-derivation-initiality} gives its
initial algebra
\[
 \mathfrak P=((P^{\mathrm{all}},c^{\mathrm{all}}),\alpha),
 \qquad c^{\mathrm{all}}:P^{\mathrm{all}}\to F_E,
\]
where
$\alpha:H_E(P^{\mathrm{all}},c^{\mathrm{all}})
             \to(P^{\mathrm{all}},c^{\mathrm{all}})$
is a morphism in $\mathcal B/F_E$.
Let $c:P\to E\times S$ be the pullback of $c^{\mathrm{all}}$
along $1_E\times j:E\times S\to F_E$, and let
$m:\mathcal T\hookrightarrow E\times S$ be the regular image
of $c$. Theorem~\ref{thm:least-theorem-subobject} identifies
the image of $c^{\mathrm{all}}$ as the least subobject closed
under these rules. Stability of regular images identifies
$m$ with the pullback of that image along $1_E\times j$.
By Lemma~\ref{lem:syntactic-total-category}, $m$ has a partial
inverse $r:E\times S\rightharpoonup\mathcal T$ in
$\Split(S_{\PA})$. Thus $rm=1_{\mathcal T}$ and $mr=\bar r$;
Lemma~\ref{lem:range-regular-image} gives
\[
 t=mr=\ran(c):E\times S\rightharpoonup E\times S
 \quad\text{in }\Split(S_{\PA}).
\]

For $e\in\N$, let $\epsilon_e:1\to E$ be the numeral
morphism. For $X\in\mathcal B$, let $!_X:X\to1$ be the
terminal morphism, and define
$k_e^X=\langle\epsilon_e!_X,1_X\rangle:X\to E\times X$.

\begin{proposition}[Fibers of the derivation family]
\label{prop:arithmetic-theory-fibers}
For every $e\in\N$, the pullback of $c^{\mathrm{all}}$ along
$k_e^{\Fm}$ is the underlying conclusion morphism of an initial
$H_{A_{\Gamma_e}}$-algebra in $\mathcal B/\Fm$.
There is an induced morphism $\iota_e:P_{\Gamma_e}\to P$
in $\mathcal B$, and the square
\[
\begin{tikzcd}
 P_{\Gamma_e}\arrow[r,"\iota_e"]
       \arrow[d,"c_{\Gamma_e}"'] & P\arrow[d,"c"]\\
 S\arrow[r,"k_e^S"'] & E\times S
\end{tikzcd}
\]
is a pullback in $\mathcal B$. Consequently
$(k_e^S)^*m$ is the sentence theory generated by $A_{\Gamma_e}$.
\end{proposition}
\begin{proof}
Put $E_{\ne e}=(\omega,[x\ne\bar e]_{\PA})\in\mathcal B$,
and let $u_e:E_{\ne e}\hookrightarrow E$ be its inclusion.
The predicates $x=\bar e$ and $x\ne\bar e$ give the
isomorphism $[\epsilon_e,u_e]:1+E_{\ne e}\to E$.
For $X\in\mathcal B$, put
$k_{\ne e}^X=u_e\times1_X:E_{\ne e}\times X\to E\times X$.
Finite extensivity gives an equivalence
\[
 \Xi_e:\mathcal B/F_E\longrightarrow
       \mathcal B/\Fm\times\mathcal B/(E_{\ne e}\times\Fm).
\]
For an object $(Q,q)$ of $\mathcal B/F_E$, define
\[
 \Xi_e(Q,q)=((Q_e,q_e),(Q_{\ne e},q_{\ne e}))
   =\bigl((k_e^{\Fm})^*(Q,q),(k_{\ne e}^{\Fm})^*(Q,q)\bigr).
\]
Thus $q_e:Q_e\to\Fm$ and
$q_{\ne e}:Q_{\ne e}\to E_{\ne e}\times\Fm$ are pullbacks
of $q$. For a slice morphism $f:(Q,q)\to(Q',q')$, represented
by $f:Q\to Q'$ with $q'f=q$, put
\[
 \Xi_e(f)=(f_e,f_{\ne e}),\qquad
 f_e:(Q_e,q_e)\to(Q'_e,q'_e),\quad
 f_{\ne e}:(Q_{\ne e},q_{\ne e})\to(Q'_{\ne e},q'_{\ne e}),
\]
where each component is the induced pullback morphism in its
respective slice category. Pullback uniqueness proves the
identity and composition laws.
A quasi-inverse
\[
 \Psi_e:\mathcal B/\Fm\times\mathcal B/(E_{\ne e}\times\Fm)
                 \longrightarrow\mathcal B/F_E
\]
sends $(x:X\to\Fm,y:Y\to E_{\ne e}\times\Fm)$ to
$[k_e^{\Fm}x,k_{\ne e}^{\Fm}y]:X+Y\to F_E$.
For $f:(X,x)\to(X',x')$ in $\mathcal B/\Fm$ and
$g:(Y,y)\to(Y',y')$ in $\mathcal B/(E_{\ne e}\times\Fm)$,
its action on morphisms is represented by
$f+g:X+Y\to X'+Y'$ in $\mathcal B$.
The decomposition $Q\cong Q_e+Q_{\ne e}$ supplies the
equivalence isomorphisms.

Write $F_{\ne e}=E_{\ne e}\times\Fm$, and let
$\pi_E:F_E\to E$ be the product projection. For $n=0,1,2$,
put $\rho_n=\pi_Eq_n^E:R_n^E\to E$, and form the pullback
\[
 R_n^{\ne e}=E_{\ne e}\times_E R_n^E,
 \qquad s_n:R_n^{\ne e}\to R_n^E,
 \qquad t_n:R_n^{\ne e}\to E_{\ne e},
 \qquad \rho_ns_n=u_et_n.
\]
Since every premise has the parameter of the conclusion, there
are unique morphisms
\[
 p_n^{\ne e}:R_n^{\ne e}\to F_{\ne e}^n,
 \qquad q_n^{\ne e}:R_n^{\ne e}\to F_{\ne e}
\]
satisfying
\[
 (k_{\ne e}^{\Fm})^n p_n^{\ne e}=p_n^Es_n,
 \qquad k_{\ne e}^{\Fm}q_n^{\ne e}=q_n^Es_n.
\]
These are the premise and conclusion morphisms for an inference
endofunctor
$H_{\ne e}:\mathcal B/F_{\ne e}\to\mathcal B/F_{\ne e}$.
Write $H_e=H_{A_{\Gamma_e}}:\mathcal B/\Fm\to\mathcal B/\Fm$
for the inference endofunctor of $A_{\Gamma_e}$.
There is a natural isomorphism
\[
 \zeta:\Xi_eH_E\xRightarrow{\cong}
                 (H_e\times H_{\ne e})\Xi_e.
\]
Its component at $(Q,q)$ consists of two slice morphisms
\[
 \begin{aligned}
 \zeta^0_{(Q,q)}&:(k_e^{\Fm})^*H_E(Q,q)\xrightarrow{\cong}
                                      H_e(Q_e,q_e),\\
 \zeta^1_{(Q,q)}&:(k_{\ne e}^{\Fm})^*H_E(Q,q)\xrightarrow{\cong}
                                      H_{\ne e}(Q_{\ne e},q_{\ne e}).
 \end{aligned}
\]
The first belongs to $\mathcal B/\Fm$ and the second to
$\mathcal B/F_{\ne e}$.
For the nullary summand of $H_E(Q,q)$, the two pullbacks of
$R_0^E$ are the corresponding axiom objects. In the first,
the nonlogical axiom predicate is $\operatorname{Ax}(\bar e,z)$.
For $n=1,2$, the equation $p_n^E=d_n(1_E\times p_n)$ makes
the parameter of every premise equal to that of the rule.
Pullback pasting therefore gives the displayed isomorphisms on
each summand of arity $n$. Taking their coproducts defines
$\zeta^0_{(Q,q)}$ and $\zeta^1_{(Q,q)}$.
For $f:(Q,q)\to(Q',q')$ in $\mathcal B/F_E$, uniqueness
of the pullback morphisms gives the naturality equation
\[
 (H_e\times H_{\ne e})(\Xi_e f)\,\zeta_{(Q,q)}
   =\zeta_{(Q',q')}\,\Xi_e(H_Ef).
\]
Both sides are morphisms from $\Xi_eH_E(Q,q)$ to
$(H_e\times H_{\ne e})\Xi_e(Q',q')$ in the product of
the two slice categories.

The equivalence $\Xi_e$ and the natural isomorphism $\zeta$
induce an equivalence
\[
 \Lambda_e:\operatorname{Alg}(H_E)\longrightarrow
             \operatorname{Alg}(H_e)\times
             \operatorname{Alg}(H_{\ne e}).
\]
An algebra $((Q,q),\delta)$, where
$\delta:H_E(Q,q)\to(Q,q)$ in $\mathcal B/F_E$, is sent
to the pair of algebras whose structure morphisms are the
components of
\[
 \Xi_e(\delta)\,\zeta_{(Q,q)}^{-1}:
 (H_e\times H_{\ne e})\Xi_e(Q,q)\longrightarrow\Xi_e(Q,q).
\]
For an algebra morphism
$a:((Q,q),\delta)\to((Q',q'),\delta')$ in
$\operatorname{Alg}(H_E)$, let
$f:(Q,q)\to(Q',q')$ be its underlying slice morphism.
Then $\Lambda_e(a)$ has underlying slice morphisms
$(f_e,f_{\ne e})$. Naturality of $\zeta$ identifies the
algebra-morphism equations in the two categories. Full faithfulness
of $\Xi_e$ therefore gives full faithfulness of $\Lambda_e$.
Transporting a pair of structure morphisms through $\Psi_e$
gives essential surjectivity.

Let $(\mathfrak P_e,\mathfrak P_{\ne e})=\Lambda_e(\mathfrak P)$.
This pair is initial in
$\operatorname{Alg}(H_e)\times\operatorname{Alg}(H_{\ne e})$.
For any $\mathfrak Q\in\operatorname{Alg}(H_e)$, there is
therefore a unique morphism
\[
 (\mathfrak P_e,\mathfrak P_{\ne e})\longrightarrow
 (\mathfrak Q,\mathfrak P_{\ne e})
\]
in that product category. Its first component is an
$H_e$-algebra morphism $\mathfrak P_e\to\mathfrak Q$.
Pairing any two such algebra morphisms with
$1_{\mathfrak P_{\ne e}}$ proves their equality. Hence
$\mathfrak P_e$ is initial and is therefore isomorphic in
$\operatorname{Alg}(H_e)$ to the initial derivation algebra for
$A_{\Gamma_e}$ from Theorem~\ref{thm:syntactic-generated-theory}.
The underlying slice object of $\mathfrak P_e$ is
$(k_e^{\Fm})^*(P^{\mathrm{all}},c^{\mathrm{all}})$.

Pull the underlying slice isomorphism of the initial algebras back along
$j:S\hookrightarrow\Fm$. Since $(1_E\times j)k_e^S=k_e^{\Fm}j$,
pullback pasting gives an isomorphism
$P_{\Gamma_e}\cong S\times_{E\times S}P$ over $S$.
Composing it with the projection to $P$ gives $\iota_e$ in the
displayed square. Regular-image stability shows that
$(k_e^S)^*m$ is the image of $c_{\Gamma_e}$.
Appendix~\ref{app:first-order-derivations} identifies that image
with the generated sentence theory.
\end{proof}

Choose a $\Sigma_1$ predicate $P_0(d)$ and an
isomorphism $u:J(P_0)\xrightarrow{\cong}P$ in $\mathcal B$.
By fullness of $J$, there is a morphism
$\gamma_c:P_0\to D$ in $\mathcal P_{\PA}$ satisfying
\[
 J(\gamma_c)=v^{-1}cu:J(P_0)\longrightarrow J(D).
\]
Choose a standard program index $a_c\in\N$ whose graph
$G_{a_c}(d,w)$ represents $\gamma_c$, as in
Lemma~\ref{lem:syntactic-functional-search}. The defining conditions
\eqref{eq:syntactic-support} and \eqref{eq:syntactic-totality}
for $\gamma_c:P_0\to D$ give
\[
 \PA\vdash\forall d\,
     \bigl(\exists w\,G_{a_c}(d,w)\leftrightarrow P_0(d)\bigr).
\]
Define $\operatorname{Prov}(e,z)
\equiv\exists d\,G_{a_c}(d,\langle e,z\rangle)$.
The product isomorphism $v$ makes this a $\Sigma_1$
predicate for $m:\mathcal T\hookrightarrow E\times S$.
For each $e\in\N$,
Proposition~\ref{prop:arithmetic-theory-fibers} identifies the
represented subobjects of $S$, giving
\[
 \PA\vdash\forall z\,
   \bigl(\operatorname{Prov}(\bar e,z)
       \leftrightarrow\operatorname{Prov}_{\Gamma_e}(z)\bigr).
\]


\subsection{Uniform independent sentences}
\label{subsec:uniform-independent-sentences}

Fix a presentation $(A,R)$ of a consistent theory with a c.e.\ axiom
set $A$ and a uniformly c.e.\ family $R$ of finitary inference rules.
In Pour-El's terminology \cite{PourEl}, the presentation is
\emph{effectively extensible} if there is a total recursive function $g:\N\to\N$ with the following
property. Whenever $e$ indexes a c.e.\ set $W_e$ of sentences in the
same language and $A\cup W_e$ is consistent under $R$, the number
$g(e)$ codes a sentence $\varphi_e$ such that
$A\cup W_e\nvdash_R\varphi_e$ and $A\cup W_e\nvdash_R\neg\varphi_e$
\cite[pp.~272, 274, Definitions~5, 6(b) and~8]{PourEl1970}.\footnote{The
conclusion of Definition~8 on p.~274 prints $\cap$ for $\cup$. On
p.~272 Pour-El specifies a sentence outside both the theorem set and
the refutation set of the extension.}
For PA we construct this assignment in
$\mathcal B$ using the Turing morphism of $S_{\PA}$.

\begin{theorem}[Uniform independent sentences]
\label{thm:uniform-independent-sentences}
Let $E=(\omega,1_\omega)$ and let $S$ be the arithmetic sentence
object in $\mathcal B=\Tot(\Split(S_{\PA}))$. For each $e\in\N$,
let $\Gamma_e$ consist of the chosen PA axioms and all arithmetic
sentences on whose codes program $e$ halts. There is a morphism
$q:E\to S$ in $\mathcal B$ such that, for every $e\in\N$ with
$\Gamma_e$ consistent, the
sentence with code $(G_{\N}q)(e)$ is neither provable nor
refutable from $\Gamma_e$.
\end{theorem}
\begin{proof}
Use the provability predicate $\operatorname{Prov}(e,z)$ of
Section~\ref{subsec:indexed-arithmetic-theories}, the morphism
$\theta:E\to S$ of Lemma~\ref{lem:arithmetic-diagonal-representation},
and the function $\Theta:\N\to\N$ defined in
Section~\ref{subsec:arithmetic-internal-comparison}.
Define the arithmetic formulas
\[
 C(e,x)\equiv\operatorname{Prov}(e,\Theta(x)),\qquad
 C'(e,x)\equiv\operatorname{Prov}(e,\operatorname{Neg}(\Theta(x))),
\]
and
\[
 R(e,x,z)\equiv
   (C(e,x)\land z=2x+1)\lor(C'(e,x)\land z=2x).
\]
Lemma~\ref{lem:syntactic-functional-search}, with input
$\langle e,x\rangle$, gives a standard index
$a_\ell\in\mathcal I(\omega\times\omega,\omega+\omega)$
such that
\[
 \begin{aligned}
 \PA&\vdash\forall e\,\forall x\,\forall z\,
   \bigl(G_{a_\ell}(\langle e,x\rangle,z)\to R(e,x,z)\bigr),\\
 \PA&\vdash\forall e\,\forall x\,
   \bigl(\exists z\,G_{a_\ell}(\langle e,x\rangle,z)
                     \leftrightarrow C(e,x)\lor C'(e,x)\bigr).
 \end{aligned}
\]
Put $\ell=[a_\ell]_{\PA}:\omega\times\omega\rightharpoonup
\omega+\omega$ in $S_{\PA}$.

Recall the Turing morphism
$\tau_{\PA}:\omega\times\omega\rightharpoonup\omega$,
the isomorphism $b:\omega\to\omega+\omega$, and
$h=b\tau_{\PA}\Delta_\omega:\omega\rightharpoonup\omega+\omega$
in $S_{\PA}$ from Section~\ref{subsec:arithmetic-internal-comparison}.
The Turing property gives a total morphism
$g:\omega\to\omega$ in $S_{\PA}$ satisfying
\[
 \tau_{\PA}(g\times1_\omega)=b^{-1}\ell:
                    \omega\times\omega\rightharpoonup\omega.
\]
Precomposing with $\langle1_\omega,g\rangle:\omega\to\omega\times\omega$
and using $(g\times1_\omega)\langle1_\omega,g\rangle
=\langle g,g\rangle=\Delta_\omega g$ gives
\[
 hg=\ell\langle1_\omega,g\rangle:
                  \omega\rightharpoonup\omega+\omega.
\]
The morphism $I(g):E\to E$ is total, so
$q=\theta I(g):E\to S$ belongs to $\mathcal B$.

Fix $e\in\N$ with $\Gamma_e$ consistent, and put
$n=F_{\N}(g)(e)\in\N$. Choose a first-order structure
$K\models\Gamma_e$ by completeness.
Proposition~\ref{prop:standard-code-soundness} and
Proposition~\ref{prop:arithmetic-theory-fibers} show that
$C(e,x)$ and $C'(e,x)$ have disjoint truth sets in $\N$.
Let $s=F_{\N}(\ell):\N\times\N\rightharpoonup\N\amalg\N$ and
let $i_0,i_1$ be the coproduct injections. For every $x\in\N$ the
graph and domain equations for $a_\ell$, together with disjointness, give
\[
 \begin{aligned}
 ((e,x),i_0(x))\in\operatorname{graph}(s)
     &\quad\Longleftrightarrow\quad \N\models C'(e,x),\\
 ((e,x),i_1(x))\in\operatorname{graph}(s)
     &\quad\Longleftrightarrow\quad \N\models C(e,x).
 \end{aligned}
\]
At $x=n$, the identity $hg=\ell\langle1_\omega,g\rangle$
identifies $F_{\N}(h)(n)$ and $s(e,n)$ as partial values.
If $C(e,n)$ holds, the second equivalence gives $n\in D_1$,
so Lemma~\ref{lem:arithmetic-diagonal-representation} gives
$C'(e,n)$. If $C'(e,n)$ holds, the first gives $n\in D_0$,
so the same lemma gives $C(e,n)$. Disjointness rules out both cases.
For the sentence $\varphi$ with code
$(G_{\N}q)(e)=\Theta(n)$, a finite derivation of $\varphi$ or
$\neg\varphi$ from $\Gamma_e$ would witness $C(e,n)$ or $C'(e,n)$,
respectively. Hence $\varphi$ is independent of $\Gamma_e$.
\end{proof}



\subsection{The numerical realization}
\label{subsec:numerical-realization}

For $X,Y\in\mathsf O$ and $e\in\mathcal I(X,Y)$, define
$\varphi_e:|X|\rightharpoonup|Y|$ by
\[
 \varphi_e(x)=y\quad\Longleftrightarrow\quad
 \N\models G_e\bigl(\gamma_X(x),\gamma_Y(y)\bigr).
\]
Let $\mathrm{Rec}$ have the terms in $\mathsf O$ as objects and
these partial functions as morphisms, with extensional equality.
The numerical indexing represents every partial recursive function,
so this is a presentation of the recursion category of partial recursive
functions on these objects.
Its restriction idempotents are partial identities on c.e.\ subsets;
the complemented ones are partial identities on recursive subsets.

Write $\mathcal C=\mathcal C_{\mathrm{rec}}$,
$\mathscr S_{\mathrm{rec}}=\Split(\mathrm{Rec})$, and
$\mathcal B_{\mathrm{rec}}=\Tot(\mathscr S_{\mathrm{rec}})$.
An object $C$ of $\mathcal C$ is a pair $(m,C_0)$, where
$m\ge0$ and $C_0\subseteq\N^m$ is computably enumerable.
Write $|C|=C_0$ for its carrier. A morphism $f:C\to D$ is
specified by a function $|f|:|C|\to|D|$ computed by a partial
recursive program on tuple codes, defined on every input from
$|C|$. Equality is equality of these functions. In computations,
$f(x)$ abbreviates $|f|(x)$. Identities and composition are
identities and composition of the represented functions.

For each $m\geq0$, choose an injective recursive tuple coding
$\kappa_m:\N^m\to\N$ with recursive image $K_m$ and recursive
inverse $\kappa_m^{-1}:K_m\to\N^m$. Take
$\kappa_0(())=0$, $K_0=\{0\}$, and $\kappa_1=1_{\N}$;
iterated numerical pairing gives such codings for $m\geq2$.
For $C=(m,C_0)$, put $\widehat C=\kappa_m(|C|)$.


\begin{proposition}\label{prop:numerical-structure}
There is an equivalence
\[
 \operatorname{Tot}(\operatorname{Split}(\mathrm{Rec}))
       \simeq\mathcal C_{\mathrm{rec}}.
\]
The category $\mathcal C_{\mathrm{rec}}$ is regular, with finite
coproducts that are disjoint and stable under pullback. Under the
displayed equivalence, every monomorphism of $\mathcal C_{\mathrm{rec}}$
corresponds to a restriction monomorphism in
$\operatorname{Split}(\mathrm{Rec})$, and
\[
 \operatorname{Par}(\mathcal C_{\mathrm{rec}},\operatorname{Mono})
       \simeq\operatorname{Split}(\mathrm{Rec}).
\]
\end{proposition}
\begin{proof}
Define a functor $Q:\mathcal C\to\mathcal B_{\mathrm{rec}}$ by
\[
 Q(C)=(\omega,1_{\widehat C}),\qquad
 Q(f)=\kappa_n |f|\kappa_m^{-1}:
     (\omega,1_{\widehat C})\longrightarrow
     (\omega,1_{\widehat D}).
\]
For $f:C\to D$ with $C=(m,C_0)$ and $D=(n,D_0)$,
the program for $Q(f)$ first semidecides membership
in $\widehat C$, then decodes its input, computes $f$, and encodes
the output. Its domain is $\widehat C$, so the displayed
split morphism is total. The inverse coding equations prove
$Q(1_C)=1_{Q(C)}$ and $Q(gf)=Q(g)Q(f)$.

For $(X,e)\in\mathcal B_{\mathrm{rec}}$, let
$D_e=\gamma_X(\operatorname{dom}e)\subseteq\N$, and define
$F:\mathcal B_{\mathrm{rec}}\to\mathcal C$ by
\[
 F(X,e)=(1,D_e),\qquad
 |F(f)|=\gamma_Y f\gamma_X^{-1}|_{D_e}:D_e\longrightarrow D_d
\]
for a total split morphism $f:(X,e)\to(Y,d)$, so that
$F(f):F(X,e)\to F(Y,d)$ belongs to $\mathcal C$.
The restriction equation $\bar f=e$ makes the displayed function defined
on all of $D_e$, and $f=dfe$ makes its values belong to $D_d$. In the composite $F(g)F(f)$ the intermediate coding and its inverse
cancel, so $F(g)F(f)=F(gf)$. The identity of $(X,e)$ restricts to
the identity function on $D_e$.
Tuple decoding gives isomorphisms
\[
 u_C:FQ(C)\xrightarrow{\cong}C
       \quad\text{in }\mathcal C,\qquad
       |u_C|=\kappa_m^{-1}|_{\widehat C}.
\]
The restriction of $\gamma_X^{-1}$ to $D_e$ gives isomorphisms
\[
 v_{(X,e)}:QF(X,e)\xrightarrow{\cong}(X,e)
       \quad\text{in }\mathcal B_{\mathrm{rec}}.
\]
For $f:C\to D$ in $\mathcal C$ and
$a:(X,e)\to(Y,d)$ in $\mathcal B_{\mathrm{rec}}$,
the coding equations give
\[
 f u_C=u_DFQ(f),\qquad
 a v_{(X,e)}=v_{(Y,d)}QF(a).
\]
Thus $u:FQ\Rightarrow1_{\mathcal C}$ and
$v:QF\Rightarrow1_{\mathcal B_{\mathrm{rec}}}$ are natural
isomorphisms, proving the first equivalence.

The singleton $(0,\{()\})$ is terminal in $\mathcal C$.
For objects of dimensions $m,n$, the product has the Cartesian product
of their carriers, represented by concatenated tuples in $\N^{m+n}$.
Pairing concatenates the two output tuples. If $f,g:C\to D$ are
parallel morphisms, their
equalizer is the object of the same dimension as $C$ with carrier
$\{x\in|C|:f(x)=g(x)\}$, included in $C$. This carrier is c.e.:
enumerate $|C|$, run the two computations in parallel, and compare
their outputs. A morphism $h:Z\to C$
with $fh=gh$ factors through the equalizer by the same program for $h$,
with its values regarded as elements of the equalizer carrier.

The object $(1,\varnothing)$ is initial. For $C=(m,C_0)$ and
$D=(n,D_0)$, choose their coproduct as the object of dimension one
with carrier
\[
 |C+D|=\{2\kappa_m(c):c\in|C|\}
        \cup\{2\kappa_n(d)+1:d\in|D|\}\subseteq\N,
\]
with the two displayed encoding functions as injections.
For $f:C\to Z$ and $g:D\to Z$, parity determines the unique
copairing: $[f,g](2\kappa_m(c))=f(c)$ and
$[f,g](2\kappa_n(d)+1)=g(d)$.
For $f:X\to C+D$ in $\mathcal C$, let $X_C,X_D$ be the
objects of the same dimension as $X$ whose carriers are the
inverse images in $|X|$ of the even and odd summands. Both are c.e.
and their inclusions give an isomorphism $X_C+X_D\cong X$:
the inverse uses the parity of $f(x)$ to choose the coproduct injection
and retains $x$. The sets of even and odd values are disjoint, so these coproducts are disjoint and
stable under pullback.

For $f:C\to D$ in $\mathcal C$, define the image object
$\operatorname{Im}(f)$ to have the dimension of $D$ and carrier
$|\operatorname{Im}(f)|=|f|(|C|)\subseteq|D|$. This carrier is
c.e. by enumerating $|C|$ and dovetailing the computations of $f$. Its image factorization
is
\[
 C\xrightarrow{e_f}\operatorname{Im}(f)
       \xrightarrow{m_f}D,
\]
where $e_f$ is surjective and $m_f$ is the inclusion.
For any surjective morphism $q:C\to D$, define a morphism
$s_q:D\to C$ by preimage search. On input $y\in|D|$, enumerate
$|C|$ and dovetail the computations of $q$ until a value
$q(x)=y$ appears, then return $x$. Surjectivity ensures termination
on $|D|$, and $qs_q=1_D$.
If $h:C\to Z$ is constant on the fibers of $q$, put
$d=hs_q:D\to Z$. For $x\in|C|$, the elements $s_q(q(x))$ and
$x$ lie in the same fiber of $q$, so $dq=h$. Surjectivity gives
uniqueness of $d$. Hence $q$ is the coequalizer of its kernel pair.
Conversely, a regular epimorphism
cannot factor through a proper image monomorphism, so regular
epimorphisms are exactly the surjections. A pullback of a surjection
is surjective on its c.e. pullback carrier. This proves regularity.

Computable singleton maps distinguish any two elements of a carrier.
Consequently a monomorphism $m:B\to C$ in $\mathcal C$ is
injective; every injective morphism is monic. Enumerating $|B|$ and
searching for a preimage computes the partial inverse of $m$ on
$|m|(|B|)$. In numerical coordinates this gives
\[
 \mu_m:Q(C)\rightharpoonup Q(B)
       \quad\text{in }\mathscr S_{\mathrm{rec}},
\]
with
\[
 \mu_mQ(m)=1_{Q(B)},\qquad
 Q(m)\mu_m=\overline{\mu_m}.
\]
The right-hand side is the partial identity on the coded image
of $m$ in $Q(C)$.

Define a functor
\[
 \widetilde Q:\operatorname{Par}(\mathcal C,\operatorname{Mono})
                 \longrightarrow\mathscr S_{\mathrm{rec}}
\]
on objects by $\widetilde Q(C)=Q(C)$ and on a monic span by
\[
 \widetilde Q[\,C\xleftarrow{m}B\xrightarrow{g}D\,]
       =Q(g)\mu_m:Q(C)\rightharpoonup Q(D).
\]
This function is defined on the coded image of $m$ and sends
the code of $m(b)$ to the code of $g(b)$ for $b\in|B|$. An isomorphism of spans therefore leaves it
unchanged. For two composable spans, their pullback selects
the pairs $(b,b')$ for which the first right leg equals the second
left leg. The same condition is the domain condition for composing
the two partial functions. Their values then agree, proving
preservation of composition. Identity spans give identities, and
the restriction span $(m,m)$ gives the partial identity on the coded image of $m$.

Define a quasi-inverse
\[
 \mathsf Q_{\mathrm{num}}:\mathscr S_{\mathrm{rec}}
       \longrightarrow\mathcal K_{\mathrm{rec}},
 \qquad
 \mathcal K_{\mathrm{rec}}
       =\operatorname{Par}(\mathcal C,\operatorname{Mono}).
\]
On objects, put $\mathsf Q_{\mathrm{num}}(X,e)=F(X,e)=(1,D_e)$.
For a partial morphism $f:(X,e)\rightharpoonup(Y,d)$, put
\[
 \widehat f=\gamma_Y f\gamma_X^{-1}:\N\rightharpoonup\N,
 \qquad U_f=\operatorname{dom}(\widehat f)\subseteq D_e.
\]
The equation $f=dfe$ gives $\widehat f(U_f)\subseteq D_d$.
Let $\iota_f:(1,U_f)\hookrightarrow(1,D_e)$ be the inclusion,
and let $g_f:(1,U_f)\to(1,D_d)$ be the morphism with underlying
function $\widehat f|_{U_f}$. Define
\[
 \mathsf Q_{\mathrm{num}}(f)=
 [\,(1,D_e)\xleftarrow{\iota_f}(1,U_f)
                      \xrightarrow{g_f}(1,D_d)\,].
\]
The carrier $U_f$ is c.e., and $g_f$ is computable and total on it,
so both legs are morphisms of $\mathcal C$.
The identity on $(X,e)$ has $U_e=D_e$ and $g_e=1_{F(X,e)}$,
so its image is the identity span.
For $g:(Y,d)\rightharpoonup(Z,c)$,
\[
 U_{gf}=\{z\in U_f:\widehat f(z)\in U_g\},\qquad
 \widehat{gf}(z)=\widehat g(\widehat f(z))\quad(z\in U_{gf}).
\]
The pullback composing the two spans is
$P_{f,g}=(2,\{(z,w)\in U_f\times U_g:\widehat f(z)=w\})$.
Projection to the first coordinate gives an isomorphism
$P_{f,g}\to(1,U_{gf})$ in $\mathcal C$. Its inverse
$(1,U_{gf})\to P_{f,g}$ sends $z$ to $(z,\widehat f(z))$.
Both functions are computable and total on their source carriers.
This isomorphism commutes with both span legs, proving
preservation of composition.
The restriction of $f$ is the partial identity on its domain,
so $\mathsf Q_{\mathrm{num}}(\bar f)$ is the span
$(\iota_f,\iota_f)$. Thus $\mathsf Q_{\mathrm{num}}$ preserves
restrictions. On total morphisms it restricts to $F$, followed by
the canonical inclusion of $\mathcal C$ into $\mathcal K_{\mathrm{rec}}$.

Regard each $u_C:FQ(C)\to C$ as a total isomorphism in
$\mathcal K_{\mathrm{rec}}$. For a partial morphism $k:C\rightharpoonup D$
represented by $C\xleftarrow{m}B\xrightarrow{h}D$, its underlying
partial function has domain $|m|(|B|)$ and value
$|h|(|m|^{-1}(x))$ at $x$ in that domain. The inverse of $|m|$
on its image is computable by preimage search. Write $k(x)$ for
the value of this represented partial function. If a second span
$C\xleftarrow{m'}B'\xrightarrow{h'}D$ has the same domain and values,
then these spans are isomorphic. The functions $|m'|^{-1}|m|$ and $|m|^{-1}|m'|$
define inverse morphisms $w:B\to B'$ and $w':B'\to B$ in
$\mathcal C$. They satisfy $m'w=m$, $h'w=h$, $mw'=m'$, and
$hw'=h'$ by equality of the represented domains and values.
Both sides of
\[
 k u_C=u_D\mathsf Q_{\mathrm{num}}\widetilde Q(k)
       :FQ(C)\rightharpoonup D
\]
are defined at $z\in\widehat C$ exactly when
$\kappa_r^{-1}(z)$ belongs to the domain of $k$, where $C=(r,C_0)$.
There they have value $k(\kappa_r^{-1}(z))$. The span criterion gives equality of these morphisms.
For $a:(X,e)\rightharpoonup(Y,d)$, both sides of
\[
 a v_{(X,e)}
   =v_{(Y,d)}\widetilde Q\mathsf Q_{\mathrm{num}}(a)
   :QF(X,e)\rightharpoonup(Y,d)
\]
are defined at $z\in D_e$ exactly when $z\in U_a$, and their
value there is $a(\gamma_X^{-1}(z))$.
The components $u_C$ and $v_{(X,e)}$ therefore give natural
isomorphisms
\[
 u:\mathsf Q_{\mathrm{num}}\widetilde Q
             \Rightarrow1_{\mathcal K_{\mathrm{rec}}},\qquad
 v:\widetilde Q\mathsf Q_{\mathrm{num}}
             \Rightarrow1_{\mathscr S_{\mathrm{rec}}}.
\]
Thus $\widetilde Q$ and $\mathsf Q_{\mathrm{num}}$ are
quasi-inverse restriction functors.
\end{proof}

\begin{proposition}\label{prop:arithmetic-theory-realization}
Fix $d\ge0$, let $V=(1,\N)$, and let
$\Sigma_k=(1,\Sigma_{k,0})$ and $\Pi_k=(1,\Pi_{k,0})$
have recursive carriers for $0\le k\le d$. Then the
syntax and list algebras of Definition~\ref{def:categorical-syntax}
exist in $\mathcal C_{\mathrm{rec}}$ and are stable under parameters.
For every axiom subobject of the sentence object, the inference
endofunctor of Theorem~\ref{thm:first-order-theory} has an initial
algebra of finite derivations.
\end{proposition}
\begin{proof}
Use the syntax codes of Section~\ref{sec:st-parameterized-syntax}.
The primitive recursion recognizing finite syntax trees, membership
in the recursive signature objects, and numerical equality decide the carriers
of objects $\Tm,\Fm\in\mathcal C_{\mathrm{rec}}$. Each constructor is a primitive-recursive
function on its displayed source. The unique root decomposition
and the strictly smaller sizes of its children are the properties
of these codes established in Proposition~\ref{prop:st-syntax-initiality}.
Let $K:\mathcal C_{\mathrm{rec}}^2\to\mathcal C_{\mathrm{rec}}^2$
be the endofunctor of Appendix~\ref{app:first-order-syntax}, with
these coefficient objects. Its constructor algebra is
$\iota:K(\Tm,\Fm)\to(\Tm,\Fm)$.

Fix $P\in\mathcal C_{\mathrm{rec}}$. In
$(\mathcal C_{\mathrm{rec}}/P)^2$, put
\[
 \mathbf T_P=((P\times\Tm,\pi_P),(P\times\Fm,\pi_P)),
 \qquad \iota_P:K_P\mathbf T_P\to\mathbf T_P,
\]
where $\iota_P$ retains the parameter and applies the constructor
in the other coordinate. Take any $K_P$-algebra
\[
 \mathbf X=((X,\xi),(Y,\upsilon)),\qquad
 \alpha=(\alpha_1,\alpha_2):K_P\mathbf X\to\mathbf X.
\]
Here $\xi:X\to P$ and $\upsilon:Y\to P$ are morphisms of
$\mathcal C_{\mathrm{rec}}$, while
\[
 \alpha_1:K_{P,1}\mathbf X\to(X,\xi),\qquad
 \alpha_2:K_{P,2}\mathbf X\to(Y,\upsilon)
\]
are morphisms in $\mathcal C_{\mathrm{rec}}/P$.
For example, the function-symbol operation at arity $k$ is defined
on tuples $(p,s,x_1,\ldots,x_k)$ with $p\in|P|$,
$s\in|\Sigma_k|$, and $x_i\in|X|$ with $\xi(x_i)=p$
for $1\le i\le k$; its value belongs to $|X|$ and lies
over $p$. Its representing program halts on every such tuple.

On input $(p,t)\in|P|\times|\Tm|$ or
$(p,\varphi)\in|P|\times|\Fm|$, parse the tree and evaluate its
nodes in children-before-parent order. At each node, apply the
representing program for the corresponding component of $\alpha$
to its label, its children's values, and the same parameter $p$.
Induction on the number of nodes proves that the child values
have the required sort and projection to $p$. They therefore lie
in the source carrier of that component of $\alpha$; its program
halts and returns a value of the required sort over $p$.
The nullary step uses the variable or nullary-symbol operation.
Thus this algorithm is partial recursive and halts on every input
in each of the two source carriers. It defines morphisms
\[
 h_T:P\times\Tm\to X,\qquad h_F:P\times\Fm\to Y
 \quad\text{in }\mathcal C_{\mathrm{rec}},
 \qquad \xi h_T=\pi_P,\quad\upsilon h_F=\pi_P.
\]
For $h=(h_T,h_F):\mathbf T_P\to\mathbf X$, decomposition at
the root gives
\[
 h\iota_P=\alpha K_P(h):K_P\mathbf T_P\to\mathbf X
 \quad\text{in }(\mathcal C_{\mathrm{rec}}/P)^2.
\]
If $h':\mathbf T_P\to\mathbf X$ satisfies that equation,
induction on root decomposition gives $h'=h$.
Consequently $(\mathbf T_P,\iota_P)$ is initial for every $P$.
Taking $P=1$ gives the syntax algebra, and the construction for
all $P$ proves stability under parameters.

Let $W\in\mathcal C_{\mathrm{rec}}$ be the object of dimension
one whose carrier is the recursive set of finite list codes over $|V|$, with constructor algebra
$\iota^L=[\nil,\cons]:1+V\times W\to W$.
Write $w_0=|\nil|(())\in|W|$ for the empty-list value.
For $P\in\mathcal C_{\mathrm{rec}}$, let
$L_P:\mathcal C_{\mathrm{rec}}/P\to\mathcal C_{\mathrm{rec}}/P$
be the slice version of $Z\mapsto1+V\times Z$ defined in
Appendix~\ref{app:categorical-hypotheses}.
Its constructor algebra has carrier $(P\times W,\pi_P)$ and
structure $\iota^L_P:L_P(P\times W,\pi_P)\to(P\times W,\pi_P)$.
For an object $(Z,z)$ of that slice, the canonical isomorphism
$\sigma_{P,Z}:(P\times V)\times_P Z\to V\times Z$
sends $((p,v),z_0)$ to $(v,z_0)$. Its inverse sends $(v,z_0)$
to $((z(z_0),v),z_0)$. For any algebra
$\beta:L_P(Z,z)\to(Z,z)$, use this isomorphism and write
$\beta_0:P\to Z$ and $\beta_1:V\times Z\to Z$ for its
two underlying operations; thus $z\beta_0=1_P$ and
$z\beta_1=z\pi_Z$.
The recursive evaluation of a finite list gives
$h_L:P\times W\to Z$ with
\[
 \begin{aligned}
 h_L(p,w_0)&=\beta_0(p),\\
 h_L(p,\cons(v,w))&=\beta_1(v,h_L(p,w)),\qquad
 zh_L=\pi_P.
 \end{aligned}
\]
The two recursion equations give $h_L\iota^L_P=\beta L_P(h_L)$
in the slice. Induction on list length proves that the computation
halts on $|P|\times|W|$ and that every algebra morphism agrees with
$h_L$.


The free-variable and substitution algorithms in
Propositions~\ref{prop:st-list-initiality} and
\ref{prop:st-substitution-realization} use these same codes.
Their primitive-recursive traversals therefore give morphisms
\[
 \FV_{\Tm}:\Tm\to W,\quad \FV_{\Fm}:\Fm\to W,\quad
 a_T:\Tm\times(V\times\Tm)\to\Tm,\quad
 b_F:\Fm\times(V\times\Tm)\to1+\Fm
\]
in $\mathcal C_{\mathrm{rec}}$.
Proposition~\ref{prop:st-substitution-realization}
proves the carrier and constructor equations in $\PA$.
They therefore hold in $\N$. Uniqueness of the algebra morphisms in the parameter slices
identifies these numerical morphisms with the categorical ones.
In particular,
$\Adm=b_F^{-1}(\Fm\hookrightarrow1+\Fm)$ is the admissibility
subobject, and $a_F:\Adm\to\Fm$ is the induced substitution
morphism. The sentence object is
$j:S\hookrightarrow\Fm$, where
$|S|=\{s\in|\Fm|:\FV_{\Fm}(s)=w_0\}$.

Fix an axiom subobject $i_A:A\hookrightarrow S$ in
$\mathcal C_{\mathrm{rec}}$. Use $A$ as the object of axiom labels. The realized
rule data of Appendix~\ref{app:first-order-derivations} are
\[
 p_n:R_n\to\Fm^n,\qquad q_n:R_n\to\Fm\quad(0\le n\le2),
\]
and define
$H_A:\mathcal C_{\mathrm{rec}}/\Fm\to
\mathcal C_{\mathrm{rec}}/\Fm$ as in that appendix.
Use the finite derivation-tree codes of
Lemma~\ref{lem:syntactic-derivation-initiality} with these rule labels.
If $R_n$ has dimension $m_n$, represent a label by a number
$r\in\widehat{R_n}$ and decode it by $\kappa_{m_n}^{-1}$
before applying $p_n$ or $q_n$. A code is valid when each numerical
label belongs to $\widehat{R_n}$ for its node's arity $n$ and the
ordered children's conclusions are the premises of its decoded rule instance.
There are finitely many nodes. Dovetailing their membership and
computation tests semidecides validity, giving an object
$P_A^{\mathrm{all}}$ of dimension one whose carrier is the
c.e.\ set of valid derivation codes. Reading the root rule and computing its conclusion
gives $c_A^{\mathrm{all}}:P_A^{\mathrm{all}}\to\Fm$.
The node constructors give an algebra
\[
 \delta_A:H_A(P_A^{\mathrm{all}},c_A^{\mathrm{all}})
           \to(P_A^{\mathrm{all}},c_A^{\mathrm{all}})
 \quad\text{in }\mathcal C_{\mathrm{rec}}/\Fm.
\]

For any object $(Q,q)$ of $\mathcal C_{\mathrm{rec}}/\Fm$ and
algebra $\beta:H_A(Q,q)\to(Q,q)$ in this slice,
evaluate a valid derivation in children-before-parent order.
Its children have been assigned values in $|Q|$ whose
images under $q$ are their conclusions. Together with the node's
decoded rule instance these values belong to the corresponding pullback
summand of $H_A(Q,q)$. The program representing $\beta$ therefore
halts at that input and returns a value in $|Q|$ with the node's
conclusion as its image under $q$. The same argument applies to
a nullary node with no child values. Induction on tree size proves
halting throughout the finite evaluation on every input in
$|P_A^{\mathrm{all}}|$. It gives a morphism
$h_A:P_A^{\mathrm{all}}\to Q$ in $\mathcal C_{\mathrm{rec}}$ with
\[
 qh_A=c_A^{\mathrm{all}},\qquad
 h_A\delta_A=\beta H_A(h_A).
\]
The second equation is in $\mathcal C_{\mathrm{rec}}/\Fm$, from
$H_A(P_A^{\mathrm{all}},c_A^{\mathrm{all}})$ to $(Q,q)$.
Root decomposition proves the equation, and induction on tree
size proves uniqueness of $h_A$.
\end{proof}

For arithmetic, choose the signature labels and numerical
syntax codes used in Section~\ref{sec:st-parameterized-syntax}.
For a c.e.\ axiom presentation $T$ containing the chosen PA sentence
axioms, use the same axiom predicate as in
Section~\ref{subsec:arithmetic-internal-comparison}, with $\Gamma=T$.
Let
$c_T^{\mathrm{all}}:\mathrm{Prf}_T^{\mathrm{all}}\to\Fm$
be the conclusion morphism for all formulas from the proposition.
Its restriction to sentence conclusions is the pullback in $\mathcal C_{\mathrm{rec}}$
\[
\begin{tikzcd}
 \mathrm{Prf}_T\arrow[r,hook]\arrow[d,"c_T"'] &
 \mathrm{Prf}_T^{\mathrm{all}}\arrow[d,"c_T^{\mathrm{all}}"]\\
 S\arrow[r,hook,"j"'] & \Fm.
\end{tikzcd}
\]
Choose this pullback as the dimension-one object with carrier
\[
 |\mathrm{Prf}_T|
  =\{d\in|\mathrm{Prf}_T^{\mathrm{all}}|:
                 c_T^{\mathrm{all}}(d)\in |S|\}.
\]
Its upper leg is inclusion and $c_T(d)=c_T^{\mathrm{all}}(d)$
in $|S|$. Theorem~\ref{thm:first-order-theory}
identifies the regular image $m_T:\mathcal T_T\hookrightarrow S$
of $c_T$ as the sentence part of the least deductively closed
formula subobject containing the axioms of $T$.

To compare with the syntactic category, write
$S^{\mathrm{syn}}\in\mathcal B_{\PA}$ and
$c_T^{\mathrm{syn}}:P_T^{\mathrm{syn}}\to S^{\mathrm{syn}}$
for the sentence object and conclusion morphism denoted by $S,c_\Gamma$
in Section~\ref{subsec:arithmetic-internal-comparison}, at $\Gamma=T$.
Let $\Fm^{\mathrm{syn}}\in\mathcal B_{\PA}$ be the formula object
with the syntax codes fixed above. Use the axiom inclusion of the
chosen sentence-code subset in both constructions. Let
$R_n^{\mathrm{num}}$ and $R_n^{\mathrm{syn}}$ be the resulting
rule objects in $\mathcal C_{\mathrm{rec}}$ and $\mathcal B_{\PA}$,
respectively, for $0\le n\le2$. Their rule morphisms are
\[
 \begin{aligned}
 p_n^{\mathrm{num}}&:R_n^{\mathrm{num}}\to\Fm^n,&
 q_n^{\mathrm{num}}&:R_n^{\mathrm{num}}\to\Fm,\\
 p_n^{\mathrm{syn}}&:R_n^{\mathrm{syn}}\to(\Fm^{\mathrm{syn}})^n,&
 q_n^{\mathrm{syn}}&:R_n^{\mathrm{syn}}\to\Fm^{\mathrm{syn}}.
 \end{aligned}
\]
Let $m_n$ be the dimension of $R_n^{\mathrm{num}}$.
The finite-product, coproduct, and pullback comparisons for these
rule constructions give bijections
\[
 \rho_n:|R_n^{\mathrm{num}}|\xrightarrow{\cong}
                  G_{\N}R_n^{\mathrm{syn}}\quad(0\le n\le2).
\]
Here $\rho_1$ sends the generalization label $(x,\varphi)$ to the
interpreted pair with entries $x,\varphi$. For modus ponens,
$\rho_2$ sends $(\varphi,\psi)$ to the pair with these entries;
the ordered premises are $\varphi,\varphi\Rightarrow\psi$ and the
conclusion is $\psi$ in both presentations.
For $n=0$, retain the coproduct tag of the axiom.
On a nonlogical axiom, retain its sentence code.
For a logical scheme, decode its parameters and encode them in the
target product. Instantiation requires that substitution succeed,
equality substitution requires success for both terms, and quantifier
distribution requires that the variable be absent from the antecedent.
The common substitution and occurrence-list computations above
preserve these tests in both directions. Reversing the product
encodings therefore gives the inverse on each summand. These
procedures are computable on the carriers. Let
$\chi_n:G_{\N}((\Fm^{\mathrm{syn}})^n)\xrightarrow{\cong}|\Fm|^n$
be the finite-product comparison, using the fixed equality
$G_{\N}\Fm^{\mathrm{syn}}=|\Fm|$. The definitions give
\[
 \chi_n(G_{\N}p_n^{\mathrm{syn}})\rho_n=|p_n^{\mathrm{num}}|,
 \qquad
 (G_{\N}q_n^{\mathrm{syn}})\rho_n=|q_n^{\mathrm{num}}|.
\]
Let $J_{\PA}:\mathcal P_{\PA}\to\mathcal B_{\PA}$ be the
predicate equivalence of Lemma~\ref{lem:syntactic-total-category}.
For each $n$, use the predicate presentation $C_n(z)$ and
isomorphism $u_n:J_{\PA}(C_n)\xrightarrow{\cong}R_n^{\mathrm{syn}}$
chosen in the derivation construction of
Lemma~\ref{lem:syntactic-derivation-initiality}. Put
\[
 C_n(\N)=\{r\in\N:\N\models C_n(r)\},\qquad
 \eta_n=G_{\N}(u_n):C_n(\N)\xrightarrow{\cong}
                                      G_{\N}R_n^{\mathrm{syn}}.
\]
Both $\eta_n$ and its inverse are computable on their carriers,
by the representing programs of $u_n$ and $u_n^{-1}$.
At a numerical derivation node of arity $n$ with numerical label $r$,
replace that label by
$\eta_n^{-1}\rho_n(\kappa_{m_n}^{-1}(r))\in C_n(\N)$ and
recursively relabel its ordered children. The inverse replaces a
syntactic label $s\in C_n(\N)$ by
$\kappa_{m_n}\rho_n^{-1}(\eta_n(s))\in\widehat{R_n^{\mathrm{num}}}$.
Induction on tree size, using the two rule
comparison equations, proves that these procedures preserve validity
and conclusions and are inverse on the valid derivation carriers.
Restricting to sentence conclusions gives a computable bijection
\[
 \rho_T:|\mathrm{Prf}_T|\xrightarrow{\cong}G_{\N}P_T^{\mathrm{syn}},
 \qquad
 (G_{\N}c_T^{\mathrm{syn}})\rho_T=|c_T|.
\]
The sentence comparison is $G_{\N}S^{\mathrm{syn}}=|S|$.
In particular, the image of $c_T$ has underlying carrier
$\{s\in|S|:\N\models\operatorname{Prov}_T(s)\}$, for the
predicate fixed in that section.

Transport these objects and total maps along $Q$.
Define
\[
 E=Q(V)=(\omega,1_\omega),\qquad
 S_{\mathrm{rec}}=Q(S),\qquad
 t_T=\ran(Q(c_T)):S_{\mathrm{rec}}\rightharpoonup S_{\mathrm{rec}}.
\]
The partial inverse of $Q(m_T)$ and
Lemma~\ref{lem:range-regular-image} identify $t_T$ with the partial
identity on the numerical image of $c_T$.
Negation and numeral substitution are the morphisms
\[
 \nu_{\mathrm{num}}:S\to S,\quad
    \nu_{\mathrm{num}}(s)=\operatorname{Neg}(s),\qquad
 \theta_{\mathrm{num}}:V\to S,\quad
    \theta_{\mathrm{num}}(n)=\Theta(n)
\]
in $\mathcal C_{\mathrm{rec}}$. Set
$\nu=Q(\nu_{\mathrm{num}}):S_{\mathrm{rec}}\to S_{\mathrm{rec}}$
and $\theta=Q(\theta_{\mathrm{num}}):E\to S_{\mathrm{rec}}$
in $\mathscr S_{\mathrm{rec}}$.
Because the syntax codes coincide, their underlying functions are
$G_{\N}(\nu^{\mathrm{syn}})$ and $G_{\N}(\theta^{\mathrm{syn}})$,
where $\nu^{\mathrm{syn}}:S^{\mathrm{syn}}\to S^{\mathrm{syn}}$ is negation and
$\theta^{\mathrm{syn}}:\omega\to S^{\mathrm{syn}}$ substitutes a numeral
in the fixed formula $\alpha(x)$ as in
Lemma~\ref{lem:arithmetic-diagonal-representation}.


\subsection{Numerical incompleteness}
\label{sec:numerical-incompleteness}
\label{subsec:numerical-incompleteness}

Use the numerical realization and comparison functors of
Section~\ref{subsec:numerical-realization}. Let
$I:\mathrm{Rec}\to\mathscr S_{\mathrm{rec}}$ be the full restriction
functor $I(X)=(X,1_X)$, $I(f)=f$. It preserves the restriction
idempotents on each object of $\mathrm{Rec}$, their order, and their
complements. Put $E=I(\omega)=Q(V)$.

The denotation
$\tau_{\mathrm{rec}}=\varphi_u:\omega\times\omega\rightharpoonup\omega$
is a Turing morphism in $\mathrm{Rec}$ by the specialization identity
of Lemma~\ref{lem:arithmetized-evaluation}. Let the isomorphism
$\zeta:\omega\to\omega+\omega$ in $\mathrm{Rec}$ send $2m$ to
the image of $m$ under the first coproduct injection and $2m+1$
to the image of $m$ under the second. The formulas $\alpha,\beta$ of
Section~\ref{subsec:arithmetic-internal-comparison} define the sets
$D_0,D_1$ of inputs on which the diagonal computation has even and
odd output, respectively. The partial identities
$1_{D_0},1_{D_1}:\omega\rightharpoonup\omega$ in $\mathrm{Rec}$ are
therefore the two diagonal restriction idempotents for
$\tau_{\mathrm{rec}}$ and $\zeta$.

The functor $J_{\mathrm{rec}}:\mathrm{Rec}\to\Pfn$ has
$J_{\mathrm{rec}}(X)=|X|$ and $J_{\mathrm{rec}}(f)=f$.
It preserves restrictions, ranges, zero morphisms, and the chosen
products and coproducts by their numerical descriptions.
In $\Pfn$, inverse images along total morphisms preserve binary joins of restriction
idempotents.
For complementary restriction idempotents $c,c':W\rightharpoonup W$
in $\mathrm{Rec}$, let $C,C'\subseteq|W|$ be their domains. For
$x\in|W|$, parallel enumeration of $C$ and $C'$ in the coding
$\gamma_W:|W|\to\N$ decides which contains $x$. The morphism
$s:W\to W+W$ in $\mathrm{Rec}$
given by
\[
 s(x)=
 \begin{cases}
  i_1(x),&x\in C,\\
  i_0(x),&x\in C'
 \end{cases}
\]
is total recursive; here $i_0,i_1:W\to W+W$ are the coproduct
injections. Its restrictions to the two domains give
$sc=i_1c$ and $sc'=i_0c'$. Since $J_{\mathrm{rec}}(1_\omega)\ne0$,
Corollary~\ref{cor:source-conditionals} gives inseparability
of $1_{D_0},1_{D_1}$ in $\mathrm{Rec}$. A recursive separator of $D_0,D_1$ would
give complementary restriction idempotents separating $1_{D_0},1_{D_1}$, so
$D_0,D_1$ are recursively inseparable. A restriction idempotent on
$E=(\omega,1_\omega)$ in $\mathscr S_{\mathrm{rec}}$ is precisely
a partial identity on a c.e.\ subset of $\N$, with meet and join
given by intersection and union. Hence any complementary pair
separating $I(1_{D_0}),I(1_{D_1})$ would separate $1_{D_0},1_{D_1}$ in $\mathrm{Rec}$.

For a c.e.\ arithmetic presentation $T$ containing the PA sentence
axioms, use $S_{\mathrm{rec}},\nu,\theta,t_T$ from
Section~\ref{subsec:numerical-realization}. The code comparison there
identifies their underlying numerical functions and domains with
the standard interpretations in
Lemma~\ref{lem:arithmetic-diagonal-representation}. For $n\in D_0$, PA proves $\alpha(\bar n)$, and for $n\in D_1$, it
proves $\neg\alpha(\bar n)$. Since $T$ contains the PA axioms, we obtain
\begin{equation}\label{eq:arithmetic-domain-inclusions}
 I(1_{D_0})\le\theta^*t_{\PA},\qquad
 I(1_{D_1})\le\theta^*\nu^*t_{\PA},\qquad t_{\PA}\le t_T.
\end{equation}
In $\mathscr S_{\mathrm{rec}}$, the first two inequalities compare restriction
idempotents on $E$ and the third compares restriction idempotents on
$S_{\mathrm{rec}}$.
For consistent $T$, its sentence theorem set and refutation set
are disjoint, so $t_T\nu^*t_T=0_{S_{\mathrm{rec}}}$.
Joins in this splitting are partial identities on unions of c.e.\ subsets,
and inverse images along total morphisms preserve these joins. Corollary~\ref{cor:domain-incompleteness}
therefore gives inseparability of $t_T,\nu^*t_T$ and
$t_T\vee\nu^*t_T\ne1_{S_{\mathrm{rec}}}$.



\subsection{Numeral instances and theory extension}

Use the numerical objects of Section~\ref{subsec:numerical-incompleteness}.
For a consistent c.e.\ extension $T$ of $\PA$, define
\[
 P_T=\{n:T\vdash\alpha(\bar n)\},\qquad
 N_T=\{n:T\vdash\neg\alpha(\bar n)\}.
\]
These are the domains of definition of the partial identities
$\theta^*t_T$ and $\theta^*\nu^*t_T$.
They are computably enumerable disjoint sets with $A\subseteq P_T$
and $B\subseteq N_T$, and they enlarge under theory extension.
Let $\operatorname{Pair}(\mathrm{Rec})$ have objects $(X,A,B)$,
where $X\in\mathsf O$ and $A,B\subseteq J_{\mathrm{rec}}(X)$
are disjoint computably enumerable subsets.
A morphism $(X,A,B)\to(Y,P,N)$ is a total morphism $f:X\to Y$
of $\mathrm{Rec}$ satisfying $J_{\mathrm{rec}}(f)(A)\subseteq P$ and
$J_{\mathrm{rec}}(f)(B)\subseteq N$. Identities and composition are inherited
from $\mathrm{Rec}$. The assignment
$T\mapsto(\omega,P_T,N_T)$, with the identity function for theory
inclusions, is a covariant functor from the poset of consistent
computably enumerable extensions of $\PA$ to
$\operatorname{Pair}(\mathrm{Rec})$.
The inclusions $A\subseteq P_T$ and $B\subseteq N_T$ give the morphism
$1_\omega:(\omega,A,B)\to(\omega,P_T,N_T)$ in this category.
A recursive separator pulls back along a total recursive function;
consequently $P_T,N_T$ are recursively inseparable.

Bernardi's numerical argument uses a precomplete equivalence
relation $R$ on $\N$: for every partial recursive function $f$
there is a total recursive function $g$ such that
$f(n)\mathord\downarrow$ implies $f(n)\,R\,g(n)$.
For such a relation, distinct equivalence classes are recursively
inseparable; consequently the same holds for disjoint nonempty
unions of classes. The relation $\equiv_{\PA}$ on indices defined
in \eqref{eq:uniform} is one of his examples
\cite[pp.~125--126]{Bernardi}.
The argument above obtains $A,B$ from
Corollary~\ref{cor:source-conditionals} and uses finite PA-proofs to
establish $A\subseteq P_T$ and $B\subseteq N_T$.

\subsection{The associated recursion categories}
\label{subsec:arithmetic-associated-categories}

Put $\mathscr C_{\PA}=\Split(S_{\PA})$, with total category
$\mathcal B_{\PA}$. In $\mathcal B_{\PA}$ take the first-order language
with variables indexed by $\omega$ and one nullary predicate for each
index in $\omega$. These predicates are its only nonlogical symbols.
Let $\Fm^{\pm}$ and $S^{\pm}$ be its formula and sentence objects, and
let $\nu^{\pm}:S^{\pm}\to S^{\pm}$ be negation. The atomic sentence
constructor $\theta^{\pm}:\omega\to S^{\pm}$ sends a generalized index
$x:W\to\omega$ to the sentence $P_x=\theta^{\pm}x:W\to S^{\pm}$
over $W$. Theorem~\ref{thm:syntactic-generated-theory} constructs
this syntax and its finite derivations in
$\mathcal B_{\PA}$.

Choose the splittings of $a_0=[\alpha]_{\PA}$ and $a_1=[\beta]_{\PA}$
with $A_0=J_{\PA}(\alpha)$ and $A_1=J_{\PA}(\beta)$.
The inclusions $m_i:A_i\hookrightarrow\omega$ retain the numerical index.
The nonlogical axioms have the presentation
\[
 j=[\theta^{\pm}m_0,\nu^{\pm}\theta^{\pm}m_1]:
 A_0+A_1\longrightarrow S^{\pm}.
\]
Its image is $\im(\theta^{\pm}m_0)\vee\im(\nu^{\pm}\theta^{\pm}m_1)$.
In the chosen syntax codes the outer constructor distinguishes
$P_n$ from $\neg P_n$, and the atom supplies its index $n$.
This decoding inverts $j$ on its image, so a nonlogical axiom
leaf carries an index in the indicated summand.
Let $t_{\pm}:T_{\pm}\hookrightarrow S^{\pm}$ be the theory
generated by $\im(j)$.
For $x:W\to\omega$, the sentence $P_x$ factors through $t_{\pm}$
if $x$ factors through $m_0$, and $\nu^{\pm}P_x$ factors through
$t_{\pm}$ if $x$ factors through $m_1$.

\begin{lemma}[The theory generated by the diagonal pair]
\label{lem:arithmetic-signed-theory}
Assume that $\PA$ is consistent. The pair $a_0,a_1$ is inseparable
in $\mathscr C_{\PA}$, and the internal theory $t_{\pm}$ is
consistent and essentially incomplete.
\end{lemma}
\begin{proof}
Suppose that complementary restriction idempotents $d,d'$ on
$\omega$ separate $a_0,a_1$, with $a_0\le d$ and $a_1\le d'$.
Let $m:C\hookrightarrow\omega$ and $m':C'\hookrightarrow\omega$
be the total inclusions of their splitting objects in $\mathscr C_{\PA}$.
In $\mathcal B_{\PA}$ the subobjects are complementary, so
$[m,m']:C+C'\to\omega$ is an isomorphism.
For the injections $i_0,i_1:\omega\to\omega+\omega$, put
\[
 \ell=[i_1m,i_0m'][m,m']^{-1}:\omega\to\omega+\omega.
\]
Thus $\ell m=i_1m$ and $\ell m'=i_0m'$.
Its source and target in the splitting are
$(\omega,1_\omega)$ and $(\omega+\omega,1_{\omega+\omega})$.
Between identity splittings the morphisms and their restrictions are
those of $S_{\PA}$, so $\ell$ is total in $S_{\PA}$.
Apply the Turing property to the total morphism
$b^{-1}\ell\pi_2:\omega\times\omega\to\omega$.
It gives a total $g:\omega\to\omega$ with
\[
 \tau_{\PA}(g\times1_\omega)=b^{-1}\ell\pi_2.
\]
Precomposing with $\langle1_\omega,g\rangle$ gives
\[
 \tau_{\PA}\Delta_\omega g=b^{-1}\ell g,
 \qquad hg=\ell g.
\]
The composite $\ell[m,m']=[i_1m,i_0m']$ sends $C$ to the
second summand and $C'$ to the first. Since coproducts are disjoint
and pullback-stable, the inverse images of the first and second
summands under $\ell$ are represented by $m'$ and $m$, respectively.
Taking inverse images along $hg=\ell g$ therefore gives
\[
 g^*a_0=g^*d',\qquad g^*a_1=g^*d.
\]
The two separating inequalities imply $g^*d=g^*d'$.
These two subobjects are complementary: their coproduct
decomposition of $\omega$ is the pullback of $[m,m']$ along the
total morphism $g$. Their meet is therefore zero and their join is
$1_\omega$. Since they are equal, $1_\omega=0_\omega$.
This contradicts
$1_{\omega}\ne0_{\omega}$, which follows from consistency of PA.

For internal consistency of $t_{\pm}$, use the numerical
finite-derivation presentation of
Theorem~\ref{thm:syntactic-generated-theory}. PA proves
$\forall n\,\neg(\alpha(n)\land\beta(n))$, as in
Lemma~\ref{lem:arithmetic-diagonal-representation}.
Use the tagged numerical presentation of the nullary rules, so that a
positive nonlogical leaf carries its predicate index.
Consider two derivations whose conclusions are a sentence
$\varphi$ and its negation. Collect into a finite list $L$ all
labels of positive nonlogical axiom leaves occurring in either
derivation. Every member of $L$ satisfies $\alpha$.
Assign a nullary predicate $P_n$ the Boolean value true exactly
when $n$ occurs in $L$. Assign every equality the value true,
interpret negation and implication by their Boolean operations,
and give $\forall v\,\psi$ the value of $\psi$.

Extracting $L$, testing membership in $L$, and evaluating a formula
are primitive recursive on codes. In PA, every positive axiom leaf
has value true by the choice of $L$. A negative leaf has label
satisfying $\beta$, hence outside $L$, so it too has value true.
Substitution preserves the value of a formula because predicates are
nullary and every equality has value true. The logical schemes
therefore have value true, and both inference rules preserve truth.
Induction in PA on the two coded derivations gives value true to
both $\varphi$ and its negation. Thus PA proves that no such pair of
derivations exists. In the predicate presentation this gives
\[
 t_{\pm}\wedge(\nu^{\pm})^*t_{\pm}=0.
\]
Corollary~\ref{cor:signed-first-order-incompleteness} gives essential
incompleteness of $t_{\pm}$ from the inseparability of $a_0,a_1$.
\end{proof}

\begin{corollary}[Arithmetic instance of the extension theorem]
\label{cor:arithmetic-associated-categories}
Assume that $\PA$ is consistent. Let
$u:U\hookrightarrow S^{\pm}$ be any consistent internal sentence
theory containing $t_{\pm}$, and write $t_{\pm}=ui$.
Then
\[
 \Sigma_i:
 \Par(\mathcal B_{\PA}/T_{\pm})
 \longrightarrow
 \Par(\mathcal B_{\PA}/U)
\]
is fully faithful, and is not essentially surjective when $i$ is
proper. Both categories are split Turing categories, and the
theory
\[
 1_U\times u:U\times U\hookrightarrow U\times S^{\pm}
\]
in the latter category is essentially incomplete. Its partial
Boolean morphism, true on the theorem domain and false on the
refutation domain, has no total extension.

Let $\mathscr R_{\PA}[U]$ be the full restriction subcategory of
$\Par(\mathcal B_{\PA}/U)$ on the objects
$U\times X\to U$, for $X$ in $S_{\PA}$.
It is a recursion category, with
$\Split(\mathscr R_{\PA}[U])\simeq\Par(\mathcal B_{\PA}/U)$.
The functor
\[
 S_{\PA}\longrightarrow\mathscr R_{\PA}[U],\qquad
 X\longmapsto(U\times X\to U),\qquad
 f\longmapsto1_U\times f
\]
is faithful and sends $\tau_{\PA}$ to $1_U\times\tau_{\PA}$,
a Turing morphism in $\mathscr R_{\PA}[U]$.
\end{corollary}
\begin{proof}
Lemma~\ref{lem:syntactic-total-category} supplies regularity,
disjoint stable finite coproducts, and the partial inverses of
monomorphisms. The syntax of Theorem~\ref{thm:syntactic-generated-theory}
is stable under parameters. Fix $P\in\mathcal B_{\PA}$ and an axiom
subobject $A\hookrightarrow P\times S^{\pm}$ over $P$. Write $F=P\times\Fm^{\pm}$
for the formula object with projection $\pi:F\to P$.
An object over $F$ in $\mathcal B_{\PA}/P$ is determined by its
morphism $q:Q\to F$, since its projection to $P$ must be $\pi q$.
Morphisms over $F$ likewise commute with these projections. Thus the
iterated slice is equivalent to $\mathcal B_{\PA}/F$.

For a rule object $\rho_n:R_n\to P$ of arity $n>0$, write
$p_{nj}:R_n\to F$ for its $j$th premise and $c_n:R_n\to F$ for
its conclusion. These morphisms are over $P$, so
$\pi p_{nj}=\rho_n=\pi c_n$. A morphism into the pullback
$R_n\times_{F^n}Q^n$ consists of $b:Z\to R_n$ and
$d_j:Z\to Q$ with $qd_j=p_{nj}b$ for every $j$. Then
$\pi qd_j=\rho_n b$, so the same cone is a cone in the slice over $P$.
For arity zero the absolute and relative summands both reduce to
the axiom object $R_0\cong(P\times L)+A$, where $L$ is the object of logical axiom
instances. These identifications preserve the conclusion and commute
with postcomposition on the $d_j$. Their coproduct identifies the two
inference endofunctors under the slice equivalence.
Lemma~\ref{lem:syntactic-derivation-initiality} therefore gives an
initial derivation algebra for $A$ in the parameter slice.

Let $v_0$ be the variable
with code $0$. The fixed sentence
$\forall v_0\,(v_0=v_0)$ is a logical theorem, so its morphism
$1\to S^{\pm}$ factors through every $u$ under consideration.
Lemma~\ref{lem:arithmetic-signed-theory} and
Corollary~\ref{cor:domain-incompleteness} give inseparability of
the theorem and refutation domains of $u$.
The conclusions follow from
Theorems~\ref{thm:theory-extension} and
\ref{thm:associated-incompleteness}.
\end{proof}

